\section{Independent proof of the compensated fold after loop rotation}
\label{sec:independent-composition-audit}

Let $B$ denote the old-label Routing245 machine, let $R$ denote its
244-state transfer-plus-one loop rotation, and let $C$ be the combined
243-state machine.  Every machine in this section has binary alphabet,
one two-way infinite tape, head moves in $\{L,R\}$, initial state $95$,
and the target-only halt label $254$.  The number of working states
excludes that halt label.  State relabelling into the canonical tables
does not change any computation.

The precise raw SHA-256 digests of the three old-label tables are
\begin{center}\scriptsize
\begin{tabular}{cl}
$B$ & \texttt{f8dbeacdad5fd89c57e3f240781324b09110be8e0ccdb838a9c710fb28d2a267}\\
$R$ & \texttt{0160caaf5d6db52eb7eb97a76cb04db4b424eb2bc0d58617818a6c7c6a1fe776}\\
$C$ & \texttt{4207fc3d1bb16f3c64c3cff0980eaa1486160c3f467aee7685740a39e3df7300}.
\end{tabular}
\end{center}
The checker pins the first two digests and checks the third table by its
complete literal transition dictionary, as well as reporting its digest.

\subsection{The exact construction and the required scope distinction}

The rotation $B\mapsto R$ changes the rows
\[
 (12,0)\mapsto(0,R,18),\qquad
 (12,1)\mapsto(1,R,17),\qquad
 (17,1)\mapsto(0,L,20)
\]
and deletes both rows of state $26$.
The compensated fold $R\mapsto C$ changes exactly
\[
\begin{array}{c|c|c}
\text{row}&R&C\\\hline
(92,1)&(1,L,24)&(1,R,66)\\
(104,1)&(1,R,122)&(1,R,92)\\
(106,1)&(1,R,122)&(1,R,92)\\
(167,1)&(1,L,92)&(1,L,151)
\end{array}
\]
and deletes both rows of state $122$.  These deleted rows are
$(122,0)\mapsto(0,R,25)$ and $(122,1)\mapsto(1,R,66)$.
No retained row outside the seven displayed patch keys changes from $B$
to $C$.  Thus the literal inventory of $C$ has $245-2=243$ working states
and $486$ rows.  Every retained state still has both binary transitions,
and $254$ is the only target without outgoing rows.

The two edits commute as operations on literal transition dictionaries:
applying the compensated fold first to $B$ and then the rotation gives
the same table $C$.  That syntactic fact does not alone prove that the
edits preserve a computation.  The argument below verifies the stronger
fact needed here: the compensated fold's complete proof premises remain
true after the rotation.

The all-tape theorem proved in this section is specifically
$R\leftrightarrow C$ from state $95$.  The independent source and runtime
proof of the rotation supplies $B\leftrightarrow R$ for the exact blank
start and the represented register computations.  No assertion that
$B$ and $R$ agree on every arbitrary binary tape is required or made.

\subsection{The local invariant survives the rotation}

For a configuration $(q,h,T)$ of $R$, with
$T:\mathbb Z\to\{0,1\}$, put
\[
\begin{split}
I(q,h,T)\quad\Longleftrightarrow\quad&
 \bigl(q\in\{68,122,167\}\ \Longrightarrow\ T(h-1)=1\bigr)\\
&{}\land
 \bigl(q=81\ \land\ T(h-1)=1\ \Longrightarrow\ T(h-2)=1\bigr).
\end{split}
\]
It holds at state $95$ for every tape and every head position.
The complete incoming-edge lists relevant to $I$ in $R$ are
\[
\begin{array}{c|l}
68&(52,1)\mapsto(1,R,68),\quad(123,1)\mapsto(1,R,68)\\
81&(68,1)\mapsto(1,R,81),\quad(123,0)\mapsto(0,R,81)\\
122&(104,1)\mapsto(1,R,122),\quad(106,1)\mapsto(1,R,122)\\
167&(145,1)\mapsto(1,R,167),\quad(148,1)\mapsto(1,R,167).
\end{array}
\]
The rotation neither changes these rows nor creates another predecessor
of these four states.  An entry into $68$, $122$, or $167$ writes the
required left-hand $1$.  An entry into $81$ from $123$ writes a left-hand
$0$, so the implication is vacuous.  An entry into $81$ from $68$ writes
the nearer $1$ and retains the next $1$ already required at $68$.
Every other target state imposes no condition.  This proves inductiveness
of $I$ under every literal transition of $R$.

For additional independent verification, the checker exhausts all Boolean
assignments to offsets $-3,-2,-1,0,1$.  These contain every cell mentioned
by either $I$ before a step or $I$ after a left/right step, including the
cell overwritten by that step.  It checks each transition whose scanned
bit and antecedent invariant match the assignment.  This is exhaustive
evaluation of a finite implication that quantifies arbitrary exterior
tape cells.  It checks $7\,752$ admissible transition--assignment pairs
for $R$ and $7\,736$ for $C$.  The proof only needs inductiveness for
$R$; the check for $C$ supplies another consistency check of the data.

\subsection{Boundaries, ordinary transitions, and the seven local identities}

Use all working labels of $R$ except $92$ and $122$ as comparison
boundaries, together with the common halt label as a terminal boundary.
Paired boundary configurations have exactly the same state, head, and
whole tape, and satisfy $I$ in the original machine $R$.

There are $244-2=242$ working boundary states, hence $484$ rows out of
them.  Exclude the four rows
\[
 E=\{(81,1),(104,1),(106,1),(167,1)\}.
\]
The remaining $480$ rows are identical in $R$ and $C$, and no such row
enters $92$ or $122$.  All $480$ identities and all their targets are
checked directly from the two supplied tables.  Every transition from a
boundary into an excluded state therefore belongs to $E$.

The exceptional rows are covered by the following seven identities.
The starting head is written as coordinate $0$; unspecified cells are
arbitrary.  The two step counts refer respectively to $R$ and $C$.
\begin{center}
\begin{tabular}{c|l|cc|c|l}
$q$&\text{fixed initial cells}&$R$&$C$&\text{endpoint }$(q',h')$
 &\text{changed cells}\\\hline
104&$T(0)=1,\ T(1)=0$&2&4&$(25,2)$&none\\
104&$T(0)=T(1)=1$&2&2&$(66,2)$&none\\
106&$T(0)=1,\ T(1)=0$&2&4&$(25,2)$&none\\
106&$T(0)=T(1)=1$&2&2&$(66,2)$&none\\
167&$T(-1)=T(0)=1$&2&2&$(24,-2)$&none\\
81&$T(-1)=0,\ T(0)=1$&2&2&$(107,-2)$&none\\
81&$T(-2)=T(-1)=T(0)=1$&6&4&$(2,-2)$&$T(-1),T(0):1\mapsto0$
\end{tabular}
\end{center}

For the first four cases the old route is
$q,122,25$ or $q,122,66$.  The new route is $q,92,66$ when the right
bit is $1$, and $q,92,107,125,25$ when it is $0$.
The latter route visits heads $1,0,1,2$ and restores the bit just read at
every step.  The first transition itself establishes the $1$ needed on
its return to coordinate $0$.

For the $167$ case the paths are $167,92,24$ and $167,151,24$;
the needed left bit comes from $I$.  For the $81$ case with a left zero,
both paths are $81,92,107$.  For the last case the old path is
\[
 81,92,24,37,55,20,2,
\]
while the new one is
\[
 81,92,66,20,2.
\]
Both erase just the two displayed $1$ cells, preserve the next left $1$,
and end at head $-2$.  The rotation's three changed rows and its deleted
state occur in none of these paths.  All actual rows traversed by the
paths are nevertheless checked afresh, rather than inferred from that
observation.

Every displayed endpoint is a boundary and satisfies $I$ vacuously.
The four cases for $104$ and $106$ exhaust their next-bit values.
The $167$ case uses its left-bit invariant.  The two cases for $81$
exhaust the possibilities allowed by $I$.  The checker additionally
tests the partition on all $52$ five-bit assignments satisfying an
exceptional entry and $I$: exactly one displayed guard matches each.
Its local execution rejects any read outside a guard's specified
window.  Thus the equalities are whole-tape identities in arbitrary
exterior contexts, not examples on tapes artificially completed by zeros.

\subsection{The all-tape halting theorem and its finite proof}

\paragraph{Theorem.}
For every $T:\mathbb Z\to\{0,1\}$ and $h\in\mathbb Z$,
\[
 \operatorname{Halt}(R,95,h,T)
 \quad\Longleftrightarrow\quad
 \operatorname{Halt}(C,95,h,T).
\]

\paragraph{Proof.}
State $95$ is a boundary and satisfies $I$.  At a paired nonterminal
boundary, choose its single common transition if the row is ordinary;
otherwise choose the unique applicable complete block from the table.
The endpoints have identical state, head, and tape.  They are again
boundaries, and the old invariant remains true because every original
literal transition preserves it.  Induction constructs a common
sequence of paired boundaries.  Intermediate visits to boundary labels
inside a selected block need not themselves be selected for comparison.

Every block has positive duration on each side.  It has at most six
steps in $R$ and at most four in $C$.  No exceptional block contains the
halt state, and every ordinary halting transition agrees.  Therefore a
halting computation on either side ends at a paired boundary on both
sides.  If the sequence of blocks is infinite, the positive durations
give unbounded elapsed time in both machines, so neither halts.
This proves the equivalence.\hfill$\square$

The direction starting from an alleged finite halt of $C$ needs no
infinite-run principle.  If $C$ halts by time $t$ while $R$ has not halted
by $6t$, the old run provides at least $t$ complete comparison blocks.
Their new counterparts consume at least $t$ steps without encountering
a halt, a contradiction.  The converse follows by constructing the
finitely many blocks of a halting run of $R$.

\paragraph{Exact clock accounting.}
For a complete matched halting run, let $A$ count the blocks starting in
$104$ or $106$ with next bit zero, and let $D$ count those starting in
$81$ with the preceding two bits $11$.  If the two halting times are
$\tau_R$ and $\tau_C$, the seven identities give
\[
 \tau_C=\tau_R+2A-2D,
 \qquad \frac{2}{3}\tau_R\leq\tau_C\leq2\tau_R.
\]
For the inequalities, each ordinary or exceptional block separately
satisfies $(2/3)\ell_R\leq\ell_C\leq2\ell_R$; sum over blocks.
These bounds apply to matched complete halting runs.  They do not claim
that every intermediate microstep has an identical counterpart.

\subsection{A reusable criterion for transplanting the fold}

The same reasoning applies after other modifications of a baseline
machine whenever the following finite premises are checked:
\begin{enumerate}
\item The initial configuration is a comparison boundary satisfying $I$,
and the required incoming lists establish inductiveness of $I$.
\item Every retained ordinary boundary row agrees between the machines
and targets another boundary or their common halt label.
\item The guarded exceptional cases exhaust the possibilities under $I$,
and their complete finite paths have equal whole-tape endpoints that are
boundaries satisfying $I$.
\item Every block has positive duration in both machines, and no
exceptional block hides a halt.
\end{enumerate}
These are sufficient semantic premises; disjoint patch locations alone
are insufficient.  In the present construction, all the premises are
rechecked directly on $R$ and $C$.  The changed role of labels $12$ and
$17$ has no effect on them, and state $26$ is absent throughout.

\subsection{Calibration and proof-assistant status}

Write $0^{\mathbb Z}$ for the blank tape.  Combining this section with
the separately established rotation and baseline calibration gives
\[
\begin{aligned}
\operatorname{Halt}(C,95,0,0^{\mathbb Z})
&\ \Longleftrightarrow\
\operatorname{Halt}(R,95,0,0^{\mathbb Z})\\
&\ \Longleftrightarrow\
\operatorname{Halt}(B,95,0,0^{\mathbb Z})\\
&\ \Longleftrightarrow\ \neg\operatorname{Con}(\mathrm{ZFC}).
\end{aligned}
\]
The source-loop proof and its runtime certificates justify the second
line for this exact blank start.  The existing calibrated proof search
justifies the last line.  Consequently the combined machine satisfies
\[
 \mathrm{ZFC}\ \vdash\
 \bigl(\neg\operatorname{Halt}(C)\ \longleftrightarrow\
       \operatorname{Con}(\mathrm{ZFC})\bigr).
\]

The present local proof consists of finite transition identities and
induction on coded finite computations.  Its blank specialization is
arithmetically formalizable in PA.  This does not supply an existing
serialized object-language PA proof of the complete ZFC calibration.

The old loop-rotation archive predates a later proof-assistant update.
At ProveIt revision
\texttt{d91a27b462cbf7fdc6537739775fd192b06fb769}, the maintained
\texttt{routing245\_binary.md} reports closed external Lean theorems for
the exact Routing245 blank machine, with no remaining runtime, startup,
or source-bisimulation hypotheses.  Its recorded audit reports eleven
axiom probes and eleven typed callers, using only
\texttt{propext}, \texttt{Classical.choice}, and \texttt{Quot.sound}.
This independent audit inspected the maintained status; it did not
replay those historical compiler/provider caches.
The current endpoint and scope are documented at
\url{https://github.com/VladimirReshetnikov/ProveIt/blob/d91a27b462cbf7fdc6537739775fd192b06fb769/Computability/BusyBeaver/Research/ZFC/routing245_binary.md}.
No new complete Lean or Rocq endpoint for the 243-state table is claimed.

The standalone checker \texttt{audit\_composition.py} imports only Python
standard-library modules and none of the inherited compiler, source,
symbolic-runtime, simulation, or search modules.  It uses explicit
exceptions for all proof-premise failures, so optimization flags cannot
disable its checks.  Its recorded result is
\texttt{PASS\_INDEPENDENT\_243\_COMPOSITION\_AUDIT}.  The checked finite
premises support the ordinary mathematical proof above; the checker is
not itself a proof-assistant kernel.
